\begin{figure*}[!t]
\begin{tcolorbox}[
    enhanced jigsaw, 
    colframe=Black!90!white, colback=white, colbacktitle=Black!10!white, coltitle=Black, fonttitle=\bfseries,
    fontupper=\small,
    fontlower=\small,
    left=1pt,          % left padding
    right=1pt,         % right padding
    top=1pt,           % top padding
    bottom=1pt         % bottom padding
]
This protocol runs between an \emph{initiator} and a \emph{responder} node. The protocol establishes a secure communication channel using a KEM-based key exchange that replaces the traditional ECIES-based RLPx handshake.

\smallskip
\begin{itemize}[leftmargin=*]
    \item \textbf{Initiator:} Long-term KEM keypair $(\mathsf{pk}_{\text{init}}, \mathsf{sk}_{\text{init}})$. Knows the responder's peer identifier $\texttt{peerID}$ from the discovery phase.
    
    \item \textbf{Responder:} Long-term KEM keypair $(\mathsf{pk}_{\text{resp}}, \mathsf{sk}_{\text{resp}})$. 
\end{itemize}

\smallskip
\noindent{\sc Handshake Phase:}
\begin{enumerate}[leftmargin=*]
\renewcommand{\labelenumi}{\arabic{enumi}.}
    \item The initiator establishes a TCP connection to the responder.
    
    \item The responder sends their public key $\mathsf{pk}_{\text{resp}}$.
    
    \item The initiator computes $id = \mathsf{Keccak512}(\texttt{ID-HASH-DOMAIN} \parallel \mathsf{pk}_{\text{resp}})$ and checks if $id = \mathsf{peedID}$. Otherwise, abort.    
    
    \item The initiator performs KEM encapsulation, $(c_{\text{init}}, k_{\text{init}}) \leftarrow \mathsf{KEM.Encap}(\mathsf{pk}_{\text{resp}})$, generates a random nonce $n_{\text{init}}\overset{\$}{\leftarrow} \{0,1\}^{256}$ and sends the \texttt{auth} message $$\texttt{auth} = pk_{\text{init}} \parallel c_{\text{init}} \parallel n_{\text{init}}.$$
    
    \item The responder receives the \texttt{auth} message, performs KEM decapsulation $k_{\text{init}} \leftarrow \mathsf{KEM.Decap}(\mathsf{sk}_{\text{resp}}, c_{\text{init}})$ and performs KEM encapsulation with the initiator's public key, $(c_{\text{resp}}, k_{\text{resp}}) \leftarrow \mathsf{KEM.Encap}(\mathsf{pk}_{\text{init}})$.
    
    \item The responder generates a random nonce $n_{\text{resp}}\overset{\$}{\leftarrow} \{0,1\}^{256}$ and sends the \texttt{ack} message $\texttt{ack} = c_{\text{resp}} \parallel n_{\text{resp}}$.
    
    \item The initiator receives the ACK message and performs KEM decapsulation $k_{\text{resp}} \leftarrow \mathsf{KEM.Decap}(\mathsf{sk}_{\text{init}}, c_{\text{resp}})$.
\end{enumerate}

\smallskip
\noindent{\sc Key Derivation Phase:}
\begin{enumerate}[leftmargin=*]
\renewcommand{\labelenumi}{\arabic{enumi}.}
    \setcounter{enumi}{10}
    \item Both parties compute the key material $k = k_{\text{init}} \parallel k_{\text{resp}}$ and the nonce hash $h = \mathsf{Keccak256}(n_{\text{init}} \parallel n_{\text{resp}})$. 
    Then, they derive the session keys as follows:
    \begin{align*}
    s_\text{shared} &= \mathsf{Keccak256}(k \parallel h) \\
    s_\text{aes}  &= \mathsf{Keccak256}(k \parallel s_\text{shared} ) \\
    s_\text{mac}  &= \mathsf{Keccak256}(k \parallel s_\text{aes} ),
    \end{align*}
    where $s_\text{shared}$, $s_\text{aes}$ and $s_\text{mac}$ are the shared secret, AES secret and MAC secret, respectively.
\end{enumerate}
\end{tcolorbox}
\vspace{-0.2in}
\caption{RLPx handshake protocol with KEM-based key exchange. The protocol uses a bilateral KEM handshake where both parties contribute shared secrets, which are then combined to derive session keys.}
\label{fig:rlpx-handshake}
\vspace{-0.1in}
\end{figure*}